\documentclass[11pt]{article}
\usepackage[margin=1in]{geometry}
\usepackage{amsmath,amssymb,amsthm,array,longtable}
\newtheorem{theorem}{Theorem}
\newcommand{\Z}{\mathbb Z}
\title{Replacement recurrence proofs for the two odd-denominator staircases}
\author{Audited integration of the 20 October base-family proofs}
\date{26 October 2026}
\begin{document}\maketitle

\paragraph{Scope and provenance.}
This replaces the recurrence arguments of Lemmas 6.13 and 6.14 in
\texttt{revised\_version2/farey\_recursion.tex}. The explicit ranks and
their symbolic case analysis were already proved in
\texttt{flattice-thm13-20261019/BASE-ANALYSIS.md}; they are not new
discoveries of this thread. That note also proves the period extensions,
covering and complete Laplacian formulas used below directly from the
base height formulas. The present proof isolates the rank argument,
retains every incoming alternative, and derives a uniform quantitative
consequence. It does not assert inheritance inside a different child field.

Fix an odd integer $d\ge3$ and write
\[
 U=(d,1),\quad V=(1-d,d+1),\quad L=\Z U+\Z V,\quad
 T=d^2+2d-1.
\]
The incoming neighbors are horizontal at even coordinate sum and vertical
at odd coordinate sum. Write $s=1+\Delta_F g$. A forbidden set contains
both inputs at an occupied $s=1$ site and at least one at an occupied
$s=0$ site. Thus it is enough to exhibit a rank decreasing to one
input when $s=1$ and to both when $s=0$.

\section*{The standard staircase}
Let
\[
 D=\{(x,y):2-d\le x\le d,\ 0\le y\le d+1,\ 1\le x+y\le d+2\}
       \cup\{(0,0),(2,d+1)\}.
\]
The prescribed tile values are
$-y(y-1)/2+\min(x-1,0)$, with an additional $1$ at the two displayed
extra points. Their compatible global extension has constants
$c_U=0,c_V=-d(d+3)/2$. On $D$, its Laplacian is $-1$ precisely at
\[
 \{x+y\text{ even},\ x=1\}
 \ \cup\
 \{x+y\text{ odd},\ 0<y<d+1,\ 1<x+y<d+2\},                 \tag{1}
\]
and is zero elsewhere. All formulas in this note are restricted to the
stated tile unless extended periodically.

Set $p=(x+y)\bmod2$. Define $\rho_D$ by the first applicable clause:
\[
\begin{array}{c|l}
\rho_D&\text{condition}\\\hline
4&(x,y)=(0,0)\text{ or }(d,1)\\
5&(x,y)=(d,0)\text{ or }(1,d+1)\\
p&y=0,\ y=d+1,\text{ or }x\le0\\
8(d-x)+p&x\ge1,\ 1\le y\le d,\ x+y\le d\\
8&x+y=d+1\\
1&x+y=d+2.
\end{array}                                                    \tag{2}
\]
These clauses cover $D$. The only nonzero differences between two tile
points congruent modulo $L$ are $\pm U,\pm V$. The $U$ overlaps have
upper endpoints $x+y=d+2$, $2\le x\le d$, and the two additional
points $(d,1),(2,d+1)$. Their paired $\rho_D$ values are respectively
$1,4,0$. The $V$ overlaps have lower endpoints $(j,0)$,
$1\le j\le d$; the paired values are $j\bmod2$ for $j<d$, and $5$
for $j=d$. Thus (2) descends to a well-defined $L$-periodic $\rho$.

\begin{theorem}\label{thm:standard}
For the standard staircase, $A(x,y)=6x+\rho(x,y)$ decreases by at
least one to both inputs when $s=0$, and to a selected input when
$s=1$. Consequently the global staircase is recurrent for every odd
$d\ge3$.
\end{theorem}
\begin{proof}
At an even $s=0$ site $(1,y)$ with $y<d$, the two rank differences
are $-6+1-8(d-1)$ and $-1$; at $(1,d)$ they are $-13,-1$.
For even $s=1$ sites, the following list is exhaustive:
\[
\begin{array}{l|c|c}
\text{center}&\text{chosen input}&A(\text{input})-A(\text{center})\\\hline
(0,0),(d,1),(2,d+1)&\text{left}&-1\\
x\le0\text{ otherwise}&\text{left}&-5\text{ or }-1\\
y=0,\ 2\le x\le d-1&\text{left}&-5\\
1\le y\le d,\ x\ge2,\ x+y\le d&\text{right}&-1\\
1\le y\le d,\ x\ge2,\ x+y=d+1&\text{right}&-1.
\end{array}                                                    \tag{3}
\]
The last row omits $(d,1)$, already treated. The only left exits in
the first two rows use representatives $(d-1,1)$ and $(1,d+1)$,
of corrections $9$ and $5$; hence the same differences apply outside
the tile.

At an odd $s=0$ site, $3\le x+y\le d$ and $1\le y\le d$.
If $x\le0$, its correction is $1$ and both vertical neighbors have
correction zero. If $x\ge1$, its correction is $8(d-x)+1$, and both
neighbors have correction at most $8(d-x)$: they are in the same
interior region, on the bottom row, or on $x+y=d+1$ with correction
$8\le8(d-x)$. The exceptional height-$4$ sites cannot be such
neighbors. Both rank differences are therefore at most $-1$.

The odd $s=1$ sites are the boundary lists in this table. Each displayed
move decreases $\rho$ by exactly one and does not change $x$:
\[
\begin{array}{l|c|l}
\text{center}&\text{move}&\text{outside input representative, if needed}\\\hline
(x,0),\ 1\le x<d&\text{down}&(x+1-d,d)\\
(d,0)&\text{up}&\\
(x,d+1),\ 2-d\le x<1&\text{down}&\\
(1,d+1)&\text{up}&(d,1)\\
x+y=1,\ 1\le y\le d-1&\text{up}&\\
x+y=d+2,\ 2\le y<d&\text{up}&(2-y,y)\\
(2,d)&\text{up}&
\end{array}                                                    \tag{4}
\]
Parity is imposed in each row. These checks establish every incoming
inequality on the tile. Since $\rho$ and $s$ are periodic and $6x$
changes by a constant under period translation, they hold globally.
A minimum-rank point of a nonempty finite forbidden set would force
another point of smaller rank in that set, a contradiction.
\end{proof}

\section*{The alternate staircase}
Let $\widehat D$ have main region
\[
 4-2d\le x\le d,\quad 0\le y\le d,\quad 2-d\le x+y\le d+2,
\]
bottom strip $\{(x,-1):1-d\le x\le0\}$, top strip
$\{(x,d+1):4-d\le x\le3\}$, and extra points
$(1-d,0),(3,d)$. Its height formula is
\[
 -\frac{y(y+1)}2+\min(0,2-x)+\min(0,d-2+x)+\max(x+y-2,0),
\]
with an additional $1$ on $x+y=2$, $1\le y\le d-1$, and at
$(1-d,0),(3,d),(1-d,-1),(3,d+1)$.
The compatible extension has $c_U=0,c_V=-d(d+3)/2-1$.

At an even site its Laplacian is $-1$ precisely on
\[
\begin{split}
 &1\le y\le d-1,\qquad x=2-d\text{ or }x=2\text{ or }x+y=2,\\
 &\text{or at }(1-d,0),(3,d).
\end{split}                                                     \tag{5}
\]
At an odd site its Laplacian is zero precisely on these six lists,
and is $-1$ elsewhere:
\[
\begin{array}{ll}
y=-1\text{ or }d+1;&y=0,x\ge1\text{ or }y=d,x\le3-d;\\
x+y=2-d,\ 1\le y\le d-2;&x+y=d+2,\ 2\le y\le d-1;\\
x+y=1,\ 0\le y\le d-2;&x+y=3,\ 2\le y\le d.
\end{array}                                                     \tag{6}
\]
Define $\widehat\rho_D$ by ordered clauses:
\[
\begin{array}{c|l}
8d&(x,y)=(1-d,0)\text{ or }(1,1)\\
8d-12&(x,y)=(3,d)\\
8y-4&x+y=2,\ 2\le y\le d-1\\
8(1-x)+p&0\le y\le d-1,\ 2-d\le x\le0,\ x+y\le1\\
8(d+1-x)+p&1\le y\le d,\ x\ge2\\
p&\text{otherwise}.
\end{array}                                                     \tag{7}
\]
Call the fourth and fifth clauses the left and right high regions.

For completeness, rank consistency is checked on the following full
overlap lists. The lower endpoints of $V$ overlaps are $(x,0)$,
$3\le x\le d$, and $(x,-1)$, $3-d\le x\le0$; both endpoint
values are the default $p$. The $U$ overlaps split into five types:
\begin{enumerate}
\item $y=-1$, $1-d\le x\le0$: both values are $p$.
\item $y=0$, $1-d\le x\le0$: the first pair has values $8d$;
the rest match because $8(1-x)=8(d+1-(x+d))$.
\item Both endpoints in the main region, $1\le y\le d-1$:
necessarily $x\le0$ and $x+y\le1$. If $x\ge2-d$, the same
left/right equality applies; otherwise both values are $p$.
\item $(3-d,d-1)\leftrightarrow(3,d)$: both values are $8d-12$.
\item $y=d$, $4-2d\le x\le3-d$: both values are $p$.
\end{enumerate}
These checks suffice for all congruent points. Indeed a fixed level
$x-dy$ has at most two representatives, differing by $U$. A class with
two levels has exactly the four representatives
\[
 (j-d,-1),\ (j,0),\ (j-2d+1,d),\ (j-d+1,d+1),\quad 3\le j\le d,
\]
whose $U,V$ adjacency graph is a rectangle. To verify the bound of two
representatives per level, a possible pair separated by $2U$ would start
at sum $-d,1-d$, or $2-d$. The first two choices give an upper
$x$ coordinate greater than $d$; the third would end at $(3,d+1)$,
whose inverse image $(3-2d,d-1)$ is outside the domain. Larger multiples
are excluded by the full coordinate-sum width. Thus (7) defines a
periodic $\widehat\rho$ on the plane.

\begin{theorem}\label{thm:alternate}
For the alternate staircase, $\widehat A(x,y)=6x+\widehat\rho(x,y)$
has the same decreasing-input property. Its global extension is
recurrent for every odd $d\ge3$.
\end{theorem}
\begin{proof}
For even $s=0$ sites, subtract $6x$ at the center. The two effective
input heights and center correction are:
\[
\begin{array}{l|c|c}
\text{center}&\text{input heights}&\widehat\rho_D\\\hline
(1-d,0),(1,1)&3,\ 8d-1&8d\\
(3,d)&8d-13,\ 7&8d-12\\
x=2-d\text{ or }2\text{ in (5)}&-5,\ 8(d-1)-1&8(d-1)\\
x+y=2,\ 2\le y\le d-1&8y-5,\ 7&8y-4.
\end{array}                                                     \tag{8}
\]
Every comparison is strict, including $d=3$ and $y=2$.
Even $s=1$ sites in either high region choose right, decreasing rank
by one; the possible right exit from $(d,1)$ is represented by $(1,0)$,
of correction $1$. All other even $s=1$ sites have correction zero
and choose left, of correction $1$, for rank drop $5$. The possible
left exit from $(4-2d,d-1)$ is represented by $(3-d,d)$.

Odd $s=0$ sites in either high region have correction $H+1$, with
both vertical neighbors of correction at most $H$. A transition from
the left region to the correction diagonal occurs only on the sum-$1$
list of (6), except its final point whose upper neighbor has correction
zero. A transition from the right region to $(3,d)$ occurs only on
an outer-diagonal list of (6). For a default-height odd $s=0$ site,
both inputs have correction zero: all transitions from default height
to a high region, the correction line, or an exceptional even point
belong to the boundary lists (6). These observations exhaust the
region interfaces and prove the two required decreases.

Odd $s=1$ sites in the left high region lie on sum $1$ and choose
down; those in the right high region lie on the outer upper diagonal
and also choose down. The correction drops by one. Default-height
odd $s=1$ sites have correction $1$ and use these height-zero inputs:
\begin{itemize}
\item Bottom strip and bottom row: down, represented respectively by
$(x+1,d)$ and $(x+1-d,d)$ if outside the tile.
\item Top strip with $x\le1$: down; its remaining odd endpoint
$(3,d+1)$: up, represented by $(2,0)$.
\item Top row with $x\le1-d$: down; its remaining odd endpoint
$(3-d,d)$: up, represented by $(2,0)$.
\item Outer lower diagonal: up inside the tile.
\item Inner sum-$3$ diagonal: up inside the tile, except the already
listed endpoint $(3-d,d)$.
\end{itemize}
These are exactly the remaining cases of (6). Since every move is
vertical, each drops the rank by one. Periodicity transports all
inequalities globally. A minimum-rank point in a finite forbidden set
gives the same contradiction as in Theorem~\ref{thm:standard}.
\end{proof}

\section*{A uniform stronger consequence}
Both proofs exclude every nonempty forbidden set whose first coordinate
is bounded below, even if it is infinite: the rank is bounded below on
such a set, whereas the forcing rules repeatedly decrease it by at
least one. For an arbitrary nonempty domain of horizontal coordinate span
$W=\sup x-\inf x<\infty$,
synchronous burning finishes in at most
\[
 6W+8d-6\quad\text{(standard)},\qquad
 6W+8d+1\quad\text{(alternate)}.                               \tag{9}
\]
Indeed $0\le\rho\le8(d-1)+1$ and $0\le\widehat\rho\le8d$.
At each retained input step the rank drops by at least one. A point
surviving $t$ rounds forces a length-$t$ sequence within the original
domain; this is impossible once $t$ exceeds its rank range. The argument
also applies to infinite vertical strips. Empty domains are immediate.

These bounds and ranks belong to the parents' actual periodic fields.
Using them on a piece of a child still requires equality on the complete
incoming stars and control of every exit. Neither this replacement nor
the old base proof establishes the final paper's hereditary column claim.
\end{document}
